\section{Universal face colimits and the remaining critical comparison}
\label{sec:colimit-critical}

The globalization theorem is exact on every compatibility cone.  The
remaining Hall problem begins when two rigid objects are incompatible.  Their
Hall product contains nonsplit extensions and may contain non-rigid middle
terms.  Rather than suppress those terms, we encode the extension problem as
a universal kernel comparison.

\subsection{The rigid-face diagram}

Fix a globalizable parameter $S\in\Qglob$.  Let
$\mathcal P_T$ be the face poset of the reachable rigid fan.  A face
$F\in\mathcal P_T$ is represented by a basic compatible rigid object
\[
 X_F=x_1\oplus\cdots\oplus x_r.
\]
Define the Hall face algebra
\begin{equation}
 \Hall_F^S
 :=\Span\left\{
       \eps_{\oplus_i x_i^{m_i}}:
       m_i\in\mathbb N
      \right\}
 \subseteq\Hall_S^\star(\MM).
 \label{eq:Hall-face-algebra}
\end{equation}
By \cref{thm:globalizable-Hall-theta}, this is a quantum affine semigroup
algebra.  If $F\subseteq F'$, direct-sum inclusion gives an injective algebra
map
\[
 \Hall_F^S\hookrightarrow\Hall_{F'}^S.
\]
Thus the face poset defines a diagram of algebras.

On the cluster side, let $\mathbb T_F^S$ be the algebra spanned by the
normalized based monomials in the cluster variables labelling $F$.  The local
Hall--theta maps give a natural isomorphism of diagrams
\begin{equation}
 \Phi_F^S:\Hall_F^S\xrightarrow{\sim}\mathbb T_F^S.
 \label{eq:diagram-isomorphism}
\end{equation}

\begin{definition}[Universal rigid-face algebra]
\label{def:universal-face-algebra}
The \emph{universal rigid-face algebra} is the algebraic colimit
\begin{equation}
 \mathscr U_T^S
 :=\colim_{F\in\mathcal P_T}\Hall_F^S
 \simeq\colim_{F\in\mathcal P_T}\mathbb T_F^S.
 \label{eq:universal-face-colimit}
\end{equation}
\end{definition}

Concretely, $\mathscr U_T^S$ is generated by the reachable rigid rays.  It
imposes every normalized quantum-torus relation between generators lying in a
common face, and imposes no extra relation merely because two generators are
incompatible.

\subsection{Two quotient maps}

All Hall face algebras are subalgebras of the one globally twisted exact Hall
algebra.  Their inclusions agree on overlaps, so the universal property gives
an algebra map
\begin{equation}
 \pi_{\mathrm H}:\mathscr U_T^S
 \longrightarrow\Hall_{\mathrm{rig}}^S,
 \label{eq:pi-H}
\end{equation}
where $\Hall_{\mathrm{rig}}^S$ denotes the subalgebra of
$\Hall_S^\star(\MM)$ generated by all reachable rigid face algebras.  This
subalgebra is allowed to contain non-rigid Hall basis elements arising from
incompatible products.  By definition, $\pi_{\mathrm H}$ is surjective.

Likewise, the face tori sit in the ambient quantum cluster skew field.  The
local theta identifications agree on overlaps, giving a surjection
\begin{equation}
 \pi_\vartheta:\mathscr U_T^S
 \longrightarrow\mathcal A_{q,T}^{\mathrm{reach},S},
 \label{eq:pi-theta}
\end{equation}
where the target is the subalgebra generated by the reachable quantum cluster
variables, equivalently by their cluster-chamber theta functions.

\begin{theorem}[Universal kernel criterion]
\label{thm:kernel-criterion}
There exists an algebra homomorphism
\begin{equation}
 \Phi_{\mathrm{glob}}^S:
 \Hall_{\mathrm{rig}}^S
 \longrightarrow
 \mathcal A_{q,T}^{\mathrm{reach},S}
 \label{eq:global-Hall-theta-map}
\end{equation}
which restricts to $\Phi_F^S$ on every rigid face algebra if and only if
\begin{equation}
 \boxed{\ker\pi_{\mathrm H}\subseteq\ker\pi_\vartheta.}
 \label{eq:kernel-inclusion}
\end{equation}
When it exists, $\Phi_{\mathrm{glob}}^S$ is unique and satisfies
\[
 \Phi_{\mathrm{glob}}^S\circ\pi_{\mathrm H}=\pi_\vartheta.
\]
It is an isomorphism if and only if the two kernels are equal.
\end{theorem}

\begin{proof}
Suppose first that $\Phi_{\mathrm{glob}}^S$ exists.  If
$u\in\ker\pi_{\mathrm H}$, then
\[
 \pi_\vartheta(u)
 =\Phi_{\mathrm{glob}}^S(\pi_{\mathrm H}(u))=0,
\]
so \cref{eq:kernel-inclusion} holds.

Conversely, assume \cref{eq:kernel-inclusion}.  Since $\pi_{\mathrm H}$ is
surjective, define
\[
 \Phi_{\mathrm{glob}}^S(\pi_{\mathrm H}(u)):=\pi_\vartheta(u).
\]
The kernel inclusion makes this independent of the chosen lift $u$.
Multiplicativity follows from multiplicativity of both quotient maps.  The
restriction to each face is the prescribed local isomorphism, and uniqueness
follows from surjectivity.  Equality of kernels is exactly the condition that
the induced quotient map be injective.
\end{proof}

\begin{remark}[The obstruction is now a relation problem]
The theorem does not claim the inclusion.  It turns the incompatible-product
question into a concrete statement about relations in one explicitly
presented algebra.  A Hall relation between words in rigid generators must be
shown to hold after replacing each generator by its theta function.  The
nonsplit extension strata determine $\ker\pi_{\mathrm H}$; scattering and
quantum theta multiplication determine $\ker\pi_\vartheta$.
\end{remark}

\subsection{Exchange elements and extension defects}

Let $x_k,x_k'$ be exchange partners across a wall and let
$M_+,M_-$ be the two normalized monomials in the common-face variables which
occur in the ordered quantum exchange relation.  In the universal face
algebra, define the exchange element
\begin{equation}
 \mathcal E_k
 :=x_kx_k'-v^{a_k}M_+-v^{b_k}M_-.
 \label{eq:exchange-element}
\end{equation}
By construction,
\[
 \mathcal E_k\in\ker\pi_\vartheta.
\]
Its Hall image
\begin{equation}
 \mathcal D_k^{\mathrm H}:=\pi_{\mathrm H}(\mathcal E_k)
 \label{eq:Hall-exchange-defect}
\end{equation}
is an exact-Hall extension defect: it is the signed combination of all
nonsplit middle terms which remains after subtracting the two cluster
monomials.  It need not vanish in the exact Hall algebra.  A global
Hall-to-theta map, if it exists, must kill every $\mathcal D_k^{\mathrm H}$.

\begin{proposition}[A necessary family of theta-zero Hall classes]
\label{prop:exchange-defects}
If \cref{eq:kernel-inclusion} holds, then every Hall exchange defect
$\mathcal D_k^{\mathrm H}$ lies in
$\ker\Phi_{\mathrm{glob}}^S$.  More generally, every relation in the theta
kernel yields a canonically defined Hall class which the global map must
annihilate.
\end{proposition}

\begin{proof}
This follows from
$\Phi_{\mathrm{glob}}^S\pi_{\mathrm H}=\pi_\vartheta$.
\end{proof}

The proposition should not be confused with the kernel criterion itself.
Elements may vanish on the theta side without vanishing on the Hall side; a
quotient map is allowed to kill them.  The true obstruction is a Hall relation
which does \emph{not} vanish after theta evaluation.

\subsection{Filtrations suggested by presentation geometry}

The orbit-normal-space theorem supplies a natural measure of non-rigidity:
\[
 \epsilon(X)=\dim\Ext^1_{\MM}(X,X).
\]
On a fixed presentation space, this is the orbit codimension.  It is
upper-semicontinuous and vanishes on the unique rigid open orbit when one
exists.  This suggests filtering incompatible Hall products by the minimum
self-extension dimension of their middle terms.  The associated degree-zero
part retains open rigid strata, while positive degrees record discriminant
contributions.

\begin{conjecture}[Extension-codimension filtration]
\label{conj:codimension-filtration}
There is a multiplicative completion or Rees construction of
$\Hall_{\mathrm{rig}}^S$ whose filtration is controlled by presentation-orbit
codimension and for which:
\begin{enumerate}[label=(\roman*),leftmargin=2.3em]
\item the degree-zero quotient is the reachable quantum theta algebra;
\item the images of the exchange defects
      \cref{eq:Hall-exchange-defect} have strictly positive filtration degree;
\item the associated graded positive-degree terms admit a critical or
      cohomological Hall interpretation.
\end{enumerate}
\end{conjecture}

The conjecture is intentionally structural.  Additivity of
$\dim\Ext^1(X,X)$ under arbitrary extensions is not asserted, so the naive
filtration by that integer alone need not be multiplicative.  A successful
construction may require stack codimension, a Harder--Narasimhan filtration,
or vanishing-cycle weights.

\subsection{Recognition of a future critical assignment}

The broad existence of mutation-covariant $3$-Calabi--Yau Hall and
Donaldson--Thomas structures is known: quiver-with-potential mutation induces
Keller--Yang equivalences, Hall scattering diagrams encode wall crossing, and
cohomological Donaldson--Thomas theory produces quantum cluster positivity
\cite{KellerYang,BridgelandScattering,Davison}.  The missing step is to attach
a canonical framed critical class to a concrete morphism
$P^-\xrightarrow fP^0$.

We can nevertheless prove that the integrated value of any satisfactory
assignment is unique.

\begin{theorem}[Mutation-recognition principle]
\label{thm:critical-recognition}
Let $(Q,W)$ be Jacobi-finite and nondegenerate, let
$A=J(Q,W)$, and fix a reachable initial seed $T$.  Suppose an assignment
\[
 f\longmapsto\mathcal I_T(f)
\]
from reachable indecomposable reduced rigid morphism objects to the ambient
quantum cluster skew field satisfies:
\begin{enumerate}[label=(\roman*),leftmargin=2.3em]
\item for the initial morphism objects $Z_{P_i}$,
      $\mathcal I_T(Z_{P_i})$ is the $i$th initial quantum cluster variable;
\item if $f$ and $f'$ are exchange partners in adjacent maximal rigid objects,
      then their assigned elements satisfy the Berenstein--Zelevinsky ordered
      quantum exchange relation with the already assigned common summands;
\item the assignment is invariant under isomorphism and compatible with the
      labelled mutation transport.
\end{enumerate}
Then for every reachable rigid morphism $f$,
\begin{equation}
 \boxed{
 \mathcal I_T(f)=\vartheta_{\Gamma_T(I(f))}=\qchar_f.}
 \label{eq:critical-recognition}
\end{equation}
For a compatible direct sum, the corresponding statement holds with the
normalized based product.
\end{theorem}

\begin{proof}
Choose a mutation path from the initial maximal rigid object to one containing
$f$.  Starting from the normalized initial values, condition (ii) determines
the new variable at each step by the quantum exchange relation.  Hence along
this path the assignment agrees with the quantum cluster pattern.  Condition
(iii) makes the result independent of the chosen labelled path.  The resulting
cluster variable is labelled by the rigid object $f$, whose cluster $g$-vector
is $\Gamma_T(I(f))$.  Davison--Mandel's chamber theorem identifies it with the
theta function at that label.  Compatible direct sums are treated by the
based toric product.
\end{proof}

\begin{corollary}[Reduction of the direct morphism-to-critical conjecture]
\label{cor:critical-reduction}
To prove that a proposed framed critical Hall class $\Xi_T(f)$ integrates to
the morphism quantum cluster variable for every reachable rigid $f$, it is
sufficient to prove:
\begin{enumerate}[label=(\roman*),leftmargin=2.3em]
\item the initial normalization for $Z_{P_i}$; and
\item one-step mutation covariance of the integrated classes.
\end{enumerate}
No separate global Laurent calculation is required.
\end{corollary}

\begin{proof}
Apply \cref{thm:critical-recognition} to the critical integration of
$\Xi_T(f)$.
\end{proof}

\begin{remark}[What remains genuinely open]
The recognition theorem proves uniqueness, not existence.  It does not
construct $\Xi_T(f)$, compare exact map-orbit strata with vanishing-cycle
cohomology, or prove the kernel inclusion
\cref{eq:kernel-inclusion}.  These are the remaining geometric problems.
\end{remark}
